\documentclass{training}
% Filigrane PAR MODULE : \uxfond (PNG A4 dilue 30 %). Cadre noir vectoriel commun.
% Proportions original 744 px : marge 20/744, trait 1/744. Exterieur BLANC.
% Cadre via eso-pic +\put (coordonnees absolues) — pas de remember picture (decalage).
\usepackage{eso-pic}

\providecommand{\uxfond}{../../resources/fond/norme_fond_paliers_30.png}

\newlength{\uxmargin}
\newlength{\uxthick}
\newlength{\uxinnerw}
\newlength{\uxinnerh}
\setlength{\uxmargin}{0.02688\paperwidth}% 20/744
\setlength{\uxthick}{0.001344\paperwidth}% 1/744
\setlength{\uxinnerw}{\dimexpr\paperwidth-2\uxmargin-2\uxthick\relax}
\setlength{\uxinnerh}{\dimexpr\paperheight-2\uxmargin-2\uxthick\relax}

% Contenu strictement a l'interieur du trait (+ jour confortable).
\geometry{%
  includeheadfoot,
  left=22mm,
  right=22mm,
  top=20mm,
  bottom=18mm,
  headheight=30pt,
  headsep=5mm,
  footskip=12mm%
}

\AddToShipoutPictureBG{%
  \AtPageLowerLeft{%
    % Fond blanc page entiere
    \put(0,0){\color{white}\rule{\paperwidth}{\paperheight}}%
    % Filigrane module a l'interieur du cadre
    \put(\LenToUnit{\dimexpr\uxmargin+\uxthick\relax},\LenToUnit{\dimexpr\uxmargin+\uxthick\relax}){%
      \includegraphics[width=\uxinnerw,height=\uxinnerh]{\uxfond}}%
    % Cadre noir : bas, haut, gauche, droite
    \put(\LenToUnit{\uxmargin},\LenToUnit{\uxmargin}){%
      \color{black}\rule{\dimexpr\paperwidth-2\uxmargin\relax}{\uxthick}}%
    \put(\LenToUnit{\uxmargin},\LenToUnit{\dimexpr\paperheight-\uxmargin-\uxthick\relax}){%
      \color{black}\rule{\dimexpr\paperwidth-2\uxmargin\relax}{\uxthick}}%
    \put(\LenToUnit{\uxmargin},\LenToUnit{\uxmargin}){%
      \color{black}\rule{\uxthick}{\dimexpr\paperheight-2\uxmargin\relax}}%
    \put(\LenToUnit{\dimexpr\paperwidth-\uxmargin-\uxthick\relax},\LenToUnit{\uxmargin}){%
      \color{black}\rule{\uxthick}{\dimexpr\paperheight-2\uxmargin\relax}}%
  }%
}

\makeatletter
\fancyhf{}
\fancyheadoffset{0pt}
\renewcommand{\headrulewidth}{0pt}
\lhead{\parbox[t]{0.56\textwidth}{\RaggedRight\footnotesize\sloppy\@title}}
\rhead{\parbox[t]{0.36\textwidth}{\RaggedLeft\footnotesize\sloppy\@subtitle}}
\cfoot{\thepage}
% Coupures : underscore et slash dans les noms de fichiers / consignes
\DeclareRobustCommand{\_}{\textunderscore\allowbreak}
\makeatother

% Evite les depassements de longues lignes (Github, Name_FirstName_...)
\sloppy
\emergencystretch=3em
\hbadness=10000

% Sommaire type 42 : chapitres seulement
\setcounter{tocdepth}{0}
\renewcommand{\cftchappresnum}{}
\renewcommand{\cftchapaftersnum}{}
\setlength{\cftchapnumwidth}{2.0em}
\setlength{\cftbeforechapskip}{0.35em}
\cftsetindents{chapter}{0em}{2.0em}


\begin{document}

\title{Going on holidays with friends}
\subtitle{Piscine UX/UI --- UX}
\member{Piscine UX/UI}{https://github.com/01-edu/UX-UI}
\summary{From an interview guide to a prototype of a shared trip.}
\version{1.0}
\prerequisite{The previous module of this piscine, or none if this is your first.}
\maketitle

\newpage
\tableofcontents
\newpage


\chapter{Instructions}

The foreword is not an exercise. It will not be graded. The exercises start at the next chapter, from the easiest contract to the hardest. A later exercise is not taken into account if an earlier one does not hold.

Only what is in the turn-in directory of an exercise is looked at. Do not add files that the exercise does not name. The name of a deliverable follows the pattern in the exercise: \texttt{Name\_FirstName\_DeliverableName\_Date\_VersionNumber}.

You will not be given the gesture. The examples in an exercise sometimes carry a constraint that the sentence does not repeat. Read them. Your other references are the manual of the tool you picked, the person sitting next to you, and the links printed in the exercise.

The order is the order of the quest. Skipping a step because a later screen looks more interesting produces a screen nobody can explain.

\subsection*{About the pictures}

The picture in the foreword is not a hint. The pale steps behind the text are the page ground.

\subsection*{About machines that write}

You may use a tool to look up a term, the way you would use a dictionary. You may not replace an interview, a test, or a file with a generated text and sign it. If you cannot point at the person, the screen, or the line you claim, the exercise is not done.


\chapter{Foreword}

Monsieur Hulot arrives at the beach with a suitcase and no product roadmap (\emph{Les Vacances de Monsieur Hulot}, Jacques Tati, 1953). The joke of the film is that everyone shares the same shore and nobody shares the same plan. Your job this quest is to find that mismatch by talking to people, not to restage Tati.

\begin{center}
\includegraphics[width=0.78\textwidth,height=0.32\textheight,keepaspectratio]{../../resources/jeunecrack.png}\\[0.4em]
{\footnotesize Picture: JeuneCrack}
\end{center}

Recommended watching. The videos are context. They are not the deliverable.
\begin{itemize}
\item UX - Introduction \& Quest n°1, 01Edu \url{https://www.youtube.com/watch?v=UbE6WMeuvQg&list=PLHyAJ_GrRtf-RAtfngbRrFh1XDTnyXFDe}
\end{itemize}

\chapter{Exercise 00: UX II - EX 1 - Interview guide}
\exnumber{00}
\extitle{UX II - EX 1 - Interview guide}
\exfiles{ex00/PDF, named Name\_FirstName\_DeliverableName\_Date\_VersionNumber}
\exauthorize{Those named in the exercise}
\makeheaderfiles

\textbf{Context:}

Welcome to the 2nd quest.

Your mission, should you choose to accept it, is to focus on going on holidays with friends. As you may have noticed, experiencing weekends or other kinds of holidays is super nice, but may lead to some moments of irritation.

So this project's goal is to make holidays more comfortable fo everyone. But first, as in all design processes, you need to understand the real problem behind stressful moments!

\textbf{Instructions:}

You need to create an interview guide to get information from your user.

And remember, you are not your user!

\textbf{Deliverable:}

\begin{itemize}
\item A PDF document with an interview guide
\item Upload the PDF document on Github, labeled as follows: "Name\_FirstName\_DeliverableName\_Date\_VersionNumber". For example, it can be labeled as "Doe\_John\_MidFiPrototype\_05242024\_V1".
\item Add a title within the document.
\item Pay attention to the global aspect of the document. It must be clear, simple and easy to read. You can get inspiration on canva.com but don't overload your design with too much details!
\end{itemize}

\textbf{Tips:}
Here are the recommended steps :

1. Define your objective, what you are looking for when interviewing people. Usually, this has to do with your user's habits, and the most painful problems they have regarding this topic.
2. Write down 1 to 3 hypotheses you have regarding this topic, that needs to be either confirmed or disproved.
3. Pick 3 to 5 main topics you need to get information on.
4. For each topic, write 1 to 3 open-ended questions. Remember, your user cannot answer with "Yes" or "No!" Ex: How, Why, What, When, Tell me about the last time you did... etc. If you really need to introduce a new topic by a close-ended question, ask why right after.
5. Write some screening questions to have a few pieces of demographics (age, job, location, etc.)
6. Structure your questions in a funnel : start with broad questions to focus on specific topics or situations. Remember to stay neutral in your tone to let the user tell their own story. Try to focus on 10 main questions.
7. Write a greeting introduction, describing briefly (1) who you are and (2) what you are looking for. No need to say too much, but enough to provide a framework and the purpose of the interview.
8. Write a conclusion to thank your user for their time, and prepare the next steps. Ex: ask to be put in touch with someone they know and who could be interested in this topic, etc.

\textbf{Resources:}

\begin{itemize}
\item Respondent - Tool to find users \url{https://www.respondent.io}
\item How to conduct user interviews \url{https://uxdesign.cc/how-to-conduct-user-interviews-fe4b8c34b0b7?sk=9625632a522a9309a91297191e668565}
\item How to design better products through user interviews \url{https://uxdesign.cc/how-to-design-better-products-through-user-interviews-4c5142bb1fc4}
\item Asking the right questions \url{https://uxdesign.cc/asking-the-right-questions-on-user-research-interviews-and-testing-427261742a67}
\item Podcast recommendation: How to Ask Great UX Research Questions \url{https://open.spotify.com/episode/5UlvTKyVObeCDxkc8DAakU?si=fjMEHCdjRXWkYhwKhxIopQ}
\end{itemize}

\textbf{Quote:}

"To ask open-ended questions is the best approach, but it's easy to get into the weeds in data analysis when every answer is a paragraph or two of prose. Users quickly tire of answering many open-ended questions, which usually require a lot of typing and explanation." Norman Nielsen Group


\chapter{Exercise 01: UX II - Ex 2 - Run interviews}
\exnumber{01}
\extitle{UX II - Ex 2 - Run interviews}
\exfiles{ex01/PDF, named Name\_FirstName\_DeliverableName\_Date\_VersionNumber}
\exauthorize{Those named in the exercise}
\makeheaderfiles

\textbf{Instructions:}

Run 6 interviews. Go in pairs, one being the interviewer and the other the note taker.
Go outside! Meet random people in cafés, at the bus, on city squares, in parks, on benches, introduce yourself and ask for 15 minutes of their time. Run your interview guide and write down anything that could be helpful, like reactions, hesitations, etc. You can use a dictaphone to make sure you're not missing anything. But in this case, let the interviewee know and ask for their permission!
You can offer them a coffee to thank them, sometimes a smile suffices!
At least 3 of your interviews must be with strangers. You'd be surprised how people are happy to help and looooove talking about themselves.

Then, write down the script on a written document, and give elements of context regarding each interview (place, time, name, etc.)

\textbf{Deliverables:}

\begin{itemize}
\item A PDF document with the scripts from all the 6 interviews.
\item Upload the PDF document on Github, labeled as follows: "Name\_FirstName\_DeliverableName\_Date\_VersionNumber". For example, it can be labeled as "Doe\_John\_MidFiPrototype\_05242024\_V1".
\item Add a title within the document.
\item Pay attention to the global aspect of the document. It must be clear, simple and easy to read. You can get inspiration on canva.com but don't overload your design with too much details!
\end{itemize}


\chapter{Exercise 02: UX II - Ex 3 - Affinity diagram \& Empathy map}
\exnumber{02}
\extitle{UX II - Ex 3 - Affinity diagram \& Empathy map}
\exfiles{ex02/zip archive, named Name\_FirstName\_DeliverableName\_Date\_VersionNumber}
\exauthorize{Those named in the exercise}
\makeheaderfiles

\textbf{Instructions:}

Open your written scripts and extract the relevant information.

1. Use paper or digital sticky notes (Post-its or Miro.com) to isolate the pieces of information your interviewees shared with you.
1 idea, quote, or piece of information = 1 sticky note.
2. When all your scripts are written/typed down on sticky notes, gather them in clusters based on themes, opinions, habits, or anything they could have in common.
3. Label the clusters.
4. Draw conclusions
5. Fill an empathy map ONLY with the content from your affinity diagram.

\textbf{Deliverables:}

\begin{itemize}
\item An Affinity Diagram (pictures of the Affinity Diagram on paper or on digital tools like Miro)
\item An Empathy Map (pictures of the Empathy Map on paper or on digital tools like Miro)
\item Upload the documents on Github, in a zip folder named "ProjectTitle\_Name\_FirstName", with all project deliverables as follows: "Name\_FirstName\_DeliverableName\_Date\_VersionNumber".  For example, the first deliverable can be named as "Doe\_John\_MidFiPrototype\_05242024\_V1".
\item Add a title within all written documents.
\item Pay attention to the global aspect of the deliverables. They must be clear, simple and easy to read. You can get inspiration on canva.com but don't overload your design with too much details!
\end{itemize}

\textbf{Resources:}

\begin{itemize}
\item What is an Affinity Diagram and how to use it \url{https://miro.com/blog/create-affinity-diagrams/}
\item Affinity Diagram and Relation Examples \url{https://www.youtube.com/watch?v=VngWHIE4k9s}
\item Updates Empathy Map Canvas by Dave Gray \url{https://medium.com/the-xplane-collection/updated-empathy-map-canvas-46df22df3c8a}
\item Empathy Map by Krisztina Szerovay \url{https://uxknowledgebase.com/empathy-map-ca037e7686b6}
\end{itemize}

\textbf{Quote:}

Remember, you are not the user, so your opinions are not absolute truth.


\chapter{Exercise 03: UX II - Ex 4 - Problem statement}
\exnumber{03}
\extitle{UX II - Ex 4 - Problem statement}
\exfiles{ex03/PDF, named Name\_FirstName\_DeliverableName\_Date\_VersionNumber}
\exauthorize{Those named in the exercise}
\makeheaderfiles

\textbf{Instructions:}

Based on your previous work and your empathy map, define a problem statement that would clarify the main problem you want to tackle.

You may use the How Might We exercise in pairs.

\textbf{Deliverable:}

\begin{itemize}
\item A PDF document with a problem statement starting with a "How Might We"
\item Upload the PDF document on Github, labeled as follows: "Name\_FirstName\_DeliverableName\_Date\_VersionNumber". For example, it can be labeled as "Doe\_John\_MidFiPrototype\_05242024\_V1".
\item Add a title within the document.
\item Pay attention to the global aspect of the document. It must be clear, simple and easy to read. You can get inspiration on canva.com but don't overload your design with too much details!
\end{itemize}

\textbf{Resources:}

\begin{itemize}
\item Stuck on a problem? Just ask how might we \url{https://relab.academy/design-thinking/stuck-on-a-problem-just-ask-how-might-we/}
\item Problem statement \url{https://www.indeed.com/career-advice/career-development/what-is-a-problem-statement}
\item Using 'How Might We' Questions to Ideative on the Right Problems \url{https://www.nngroup.com/articles/how-might-we-questions/}
\end{itemize}

\textbf{Quote:}

'Out of clutter, find simplicity. From discord, find harmony. In the middle of difficulty, lies opportunity.' Albert Einstein


\chapter{Exercise 04: UX II - Ex 5 - Ideation}
\exnumber{04}
\extitle{UX II - Ex 5 - Ideation}
\exfiles{ex04/zip archive, named Name\_FirstName\_DeliverableName\_Date\_VersionNumber}
\exauthorize{Those named in the exercise}
\makeheaderfiles

\textbf{Instructions:}

To find creative solutions to your problem statement, use one of these techniques in groups of 5 people, twice.

Gather in groups of 2. During the first run, one is the facilitator and the other is the time-keeper and in charge of logistics. Switch roles for the second run.

Get as many ideas as possible within a defined timeframe, and narrow them down to 3 main ideas with a voting process (in pair or in group). Take pictures!

\textbf{Deliverables:}

\begin{itemize}
\item A PDF document that includes the script with all the steps + logistics
\item Pictures of the 2 sessions
\item A PDF document with 3 ideas that came out the 2 workshops.
\item Upload the documents on Github, in a zip folder named "ProjectTitle\_Name\_FirstName", with all project deliverables as follows: "Name\_FirstName\_DeliverableName\_Date\_VersionNumber".  For example, the first deliverable can be named as "Doe\_John\_MidFiPrototype\_05242024\_V1".
\item Add a title within all written documents.
\item Pay attention to the global aspect of the deliverables. They must be clear, simple and easy to read. You can get inspiration on canva.com but don't overload your design with too much details!
\end{itemize}

\textbf{Tips:}

Ideation techniques:

\begin{itemize}
\item Worst Idea \url{https://www.interaction-design.org/literature/article/learn-how-to-use-the-best-ideation-methods-worst-possible-idea}
\item Brainwriting \url{https://business.tutsplus.com/tutorials/how-to-use-brainwriting-for-rapid-idea-generation--cms-26451}
\item Mash Up \url{https://cdn.shopify.com/s/files/1/0259/7876/5396/files/IDEO_U_Ideation_Method_Mash-Up.pdf?33258}
\item As to the script, you can refer to "Athlete, keep hydrated" Quest, Ideation Exercise.
\end{itemize}

\textbf{Resources:}

\begin{itemize}
\item Where Good ideas come from by Steven Johnson \url{https://www.ted.com/talks/steven_johnson_where_good_ideas_come_from}
\item Creative Confidence by Tom Keller \url{https://booksvooks.com/nonscrolablepdf/creative-confidence-unleashing-the-creative-potential-within-us-all-pdf-1.html?page=1}
\item Podcast What is wrong with UX? \url{https://www.usersknow.com/podcast/2016/10/10/ua1wdsszhx2pemugxc1e0qj0a3l9cs}
\item Brainwriting \#2 \url{https://dux.typepad.com/dux/2011/01/method-2-of-100-brainwriting-brainwriting-is-an-ideation-method-for-quickly-generating-ideas-by-asking-people-to-write-thei.html}
\item Voting process at the end of an ideation workshop \url{https://www.interaction-design.org/literature/article/how-to-select-the-best-idea-by-the-end-of-an-ideation-session}
\end{itemize}

Quote:

\begin{itemize}
\item 'Brainstorming means using the brain to storm a creative problem. Do so in commando fashion, with each stormer attacking the same objective.' Alex Faickney Osborn
\end{itemize}


\chapter{Exercise 05: UX II - Ex 6 - User flow}
\exnumber{05}
\extitle{UX II - Ex 6 - User flow}
\exfiles{ex05/zip archive, named Name\_FirstName\_DeliverableName\_Date\_VersionNumber}
\exauthorize{Those named in the exercise}
\makeheaderfiles

\textbf{Instructions:}

Structure your main ideas into a comprehensive user flow.
You can include one, two, of the three main ideas from your ideation process, depending on how consistent it would be for the users.

\textbf{Deliverable:}

\begin{itemize}
\item A comprehensive user flow on a board (digital or pictures of a paper version)
\item Upload the documents on Github, in a zip folder named "ProjectTitle\_Name\_FirstName", with all project deliverables as follows: "Name\_FirstName\_DeliverableName\_Date\_VersionNumber".  For example, the first deliverable can be named as "Doe\_John\_MidFiPrototype\_05242024\_V1".
\item Add a title within all written documents.
\item Pay attention to the global aspect of the deliverables. They must be clear, simple and easy to read. You can get inspiration on canva.com but don't overload your design with too much details!
\end{itemize}

\textbf{Resources:}

\begin{itemize}
\item User flow glossary \url{https://www.productplan.com/glossary/user-flow/}
\item Site flows vs User Flows \url{https://uxmovement.com/wireframes/site-flows-vs-user-flows-when-to-use-which/}
\item How to make a User Flow diagram \url{https://www.lucidchart.com/blog/how-to-make-a-user-flow-diagram}
\end{itemize}

\textbf{Quote:}

'Each use case is represented as a sequence of simple steps, beginning with a user's goal and ending when that goal is fulfilled.' Usability.gov


\chapter{Exercise 06: UX II - Ex 7 - Prototyping}
\exnumber{06}
\extitle{UX II - Ex 7 - Prototyping}
\exfiles{ex06/Figma file, named Name\_FirstName\_DeliverableName\_Date\_VersionNumber}
\exauthorize{Those named in the exercise}
\makeheaderfiles

\textbf{Instructions:}

Give another dimension to your user flow and turn your user flow into a paper prototype (12 screens max).

\begin{itemize}
\item First, draw low-fidelity prototypes on paper. One main screen = 1 mobile frame. No micro-interactions nor multistate are needed for this exercise.
\item Test the prototype with 3 people. To do so, prepare a usability testing guide following the Design Sprint Five-Act interview \url{https://www.youtube.com/watch?v=U9ZG19XTbd4}
\item Then, draw your mid-fidelity screens on Figma with simple shapes, in black and white.
\end{itemize}

\textbf{Deliverables:}

\begin{itemize}
\item Pictures of a paper prototype
\item Wireframes (mid-fi) to upload on the Figma file
\item Upload the documents on Github, in a zip folder named "ProjectTitle\_Name\_FirstName", with all project deliverables as follows: "Name\_FirstName\_DeliverableName\_Date\_VersionNumber".  For example, the first deliverable can be named as "Doe\_John\_MidFiPrototype\_05242024\_V1".
\end{itemize}

\textbf{Tips:} Vocabulary

\begin{itemize}
\item Low-Fidelity: Paper Prototypes - Example here \url{https://miro.medium.com/max/1400/1*5nUfqqA2gjdAYHagjbHA5w.jpeg}
\item Mid-Fidelity: Wireframes - black and white, only the main elements, on digital prototype tools - Example here \url{https://miro.medium.com/max/1400/1*eGSfaxwYbxUFomYbyO6GWw.png}
\item High-Fidelity: Wireframes with color, styles, graphical details, and micro-interactions - Example here \url{https://miro.medium.com/max/1400/1*Xn0HSKAvhr4TZzC9lN5udw.gif}
\end{itemize}

\textbf{Resources:}

\begin{itemize}
\item Google Ventures - Design Sprint: The Five-Act interview
\item Example: Paper Prototyping - The Practical Beginners guide \url{https://www.uxpin.com/studio/blog/paper-prototyping-the-practical-beginners-guide/}
\item Example: Wellness App \url{https://medium.com/@ebtaskin/story-of-my-very-first-ux-ui-design-project-at-ironhack-bootcamp-955addeb760e}
\item Android Mobile Template \url{https://i.pinimg.com/originals/43/8e/de/438ede24264eb9552b635d34142eac29.png}
\item iPhone Template \url{https://www.cssauthor.com/wp-content/uploads/2014/01/iPhone-5S-Template-PSD.jpg}
\item Figma
\end{itemize}

\textbf{Quote:}

Any design is like a joke ...if you have to explain it, it's bad.


\end{document}
